\answer{OR22}
令 $V$ 为开始免费投掷前的最优价值。若当前点数为 $j$，接受的价值为 $j$；重新开始需先付1元，继续价值为 $V-1$。候选贝尔曼方程为
\[
 V=\frac16\sum_{j=1}^6\max\{j,V-1\}.
\]
$V=4$ 满足该式，因为
\[
 \frac{3+3+3+4+5+6}{6}=4.
\]
下面验证它确为全体策略的上界。令 $V_m$ 为最多投掷 $m$ 次的价值。$V_1=7/2\le4$，并且
\[
 V_{m+1}=\frac16\sum_{j=1}^6\max\{j,V_m-1\}.
\]
右边关于 $V_m$ 单调，故归纳得到 $V_m\le4$。若某策略在第 $N$ 次投掷停止且 $\E N<\infty$，强制它至迟在第 $m$ 次停止；所得趋于原所得，点数有界且费用被 $N$ 控制，故由支配收敛，其期望仍不超过4。若 $\E N=\infty$，由于最后点数至多6，期望净所得为负无穷，不可能更优。

因此最优决策为：1、2点重掷；4、5、6点接受；3点时接受或重掷均可。例如采用“首次出现至少3点即停”，最终点数在3、4、5、6上均匀，均值为 $9/2$；投掷次数为成功概率 $2/3$ 的几何变量，平均重掷费用为 $3/2-1=1/2$。于是
\[
 \boxed{\text{最大期望所得}=9/2-1/2=4}.
\]
这条策略达到上界且几乎必然结束，因而最优。
